% 《大学物理简明教程》5版，第9章解答题9.3.1—9.3.6
% 用户扫描原书/PDF第239页，高分辨率核对几何、方向和数值；
% 独立研究性解析，所有题仍待专家二校。

\subsection{解答题：感生、动生电动势与穿越磁场边界}

\sourceproblem{9.3.1}{239}{239}
\problemtext 匀强磁场$B=0.4\,\mathrm T$，
圆形闭合回路垂直磁场，面积随时间变化为
$S(t)=5t^2+3$（单位$\mathrm{cm^2}$、$t$单位秒）。
求$t=2$ s时的感应电动势大小。
\workings
穿过回路的磁通量为
$\Phi=BS=0.4(5t^2+3)\times10^{-4}\,\mathrm{Wb}$。
法拉第定律给
\[
|\mathcal E|=\left|\frac{d\Phi}{dt}\right|
=B\,\frac{dS}{dt}
=0.4(10t)\times10^{-4}\,\mathrm V.
\]
$t=2$ s时
\[
\boxed{|\mathcal E|=8.0\times10^{-4}\,\mathrm V}.
\]
\crosscheck $1\,\mathrm{cm^2}=10^{-4}\,\mathrm{m^2}$；
若遗漏面积单位换算将导致$10^4$倍错误。

\sourceproblem{9.3.2}{239}{239}
\problemtext 长直导线载流$I$、电流按图向下，
其右侧同平面有半圆形金属杆$MeN$，半径$b$，
圆心$O$距直导线$a$且$a>b$。
半圆杆整体以速率$v$沿直导线向上匀速移动。
求杆中动生电动势的大小和端点电势差$U_M-U_N$。
\workings
设从长直导线向右量$x$，
杆端横坐标$x_M=a-b,x_N=a+b$。
由于直导线向下通电，
杆位置处磁场垂直纸面\emph{向外}，大小
\[
B(x)=\frac{\mu_0I}{2\pi x}.
\]
速度向上，故$\bm v\times\bm B$沿$x$正向，
动生电动势由下式积分：
\[
\mathcal E_{M\to N}
=\int_M^N(\bm v\times\bm B)\cdot d\bm l
=v\int_{a-b}^{a+b}B(x)\,dx
=\boxed{\frac{\mu_0Iv}{2\pi}
\ln\frac{a+b}{a-b}}.
\]
导体上自由电荷达到准静态平衡时，
$\bm E_{\rm static}=-\bm v\times\bm B$，
右端$N$具有较高静电势，因此
\[
\boxed{U_M-U_N
=-\frac{\mu_0Iv}{2\pi}\ln\frac{a+b}{a-b}}.
\]
\crosscheck 结果与半圆弧具体形状无关，
只与杆两端的$x$坐标有关，因为$vB(x)$仅为$x$的函数；
若反向电流或反向速度，则端点电势差换号。

\sourceproblem{9.3.3}{239}{239}
\problemtext 两根平行无限长直导线在同一平面内，
电流方向相反、大小均为$I(t)$，且$\dot I>0$。
在两导线右侧放置矩形线圈，线圈高$\ell$、宽$a$；
从靠右导线到线圈近边的距离为$d$，
从靠左导线到线圈近边距离为$b$（$b>d$）。
求矩形中磁通量及感应电动势。
\workings
设右边导线的磁场\emph{指向纸内}为正通量，
两根相反电流产生的场分别为
$\mu_0I/(2\pi x)$和反向的场。
对线圈宽度逐条积分，
\[
\Phi_B(t)=\frac{\mu_0I(t)\ell}{2\pi}
\left[
\int_d^{d+a}\frac{dx}{x}
-\int_b^{b+a}\frac{dx}{x}
\right].
\]
因此
\[
\boxed{\Phi_B=
\frac{\mu_0I\ell}{2\pi}
\ln\frac{(d+a)b}{d(b+a)}},
\]
\[
\boxed{\mathcal E_{\rm ccw}=
-\frac{\mu_0\ell\dot I}{2\pi}
\ln\frac{(d+a)b}{d(b+a)}}.
\]
这里“ccw”表示按纸面向外为正法向的有向积分定义。
\crosscheck 对$b>d$，$\ln[(d+a)b/(d(b+a))]>0$，
说明近导线的贡献较强；
$\dot I>0$时感应电流反抗磁通量增大。

\sourceproblem{9.3.4}{239}{239}
\problemtext 长直导线电流$I=5$ A向上，
右侧同平面矩形线圈高$b=0.06$ m、宽$a=0.04$ m，
近边距导线$d=0.05$ m。
线圈以$v=0.03\,\mathrm{m/s}$向右远离长直导线。
求该瞬间感应电动势的大小和方向。
\workings
导线右侧的磁场指向纸内，
其通过矩形回路的磁通量大小为
\[
\Phi(d)=\frac{\mu_0Ib}{2\pi}
\int_d^{d+a}\frac{dx}{x}
=\frac{\mu_0Ib}{2\pi}\ln\frac{d+a}{d}.
\]
由于$d'(t)=v>0$，
\[
\left|\frac{d\Phi}{dt}\right|
=\frac{\mu_0Ibv}{2\pi}
\left(\frac1d-\frac1{d+a}\right).
\]
代入题设数据，
\[
\boxed{|\mathcal E|
=2.0\times10^{-7}\cdot5\cdot0.06\cdot0.03
\left(20-\frac1{0.09}\right)
=1.60\times10^{-8}\,\mathrm V}.
\]
线圈离导线越来越远，指向纸内的磁通量减小，
为阻止这种变化，感应电流产生\emph{纸内}磁场，
故\textbf{从纸外观察电流为顺时针}。
\crosscheck 面积的宽$a$与高度$b$不能互换：
靠近导线的两个长边$B$不同，
应沿径向宽度积分。

\sourceproblem{9.3.5}{239}{239}
\problemtext 导轨上导体杆$ab$长$\ell$，
以恒定速度$v$向右运动；
磁场$B(t)=kt$（$k>0$），磁感应方向与
导轨平面法线夹角为$60^\circ$。
$t=0$时杆位于固定端$cd$，
求任意时刻闭合回路的感应电动势大小和绕向。
\workings
平面有效法向磁场
$B_\perp(t)=kt\cos60^\circ=kt/2$。
杆从固定端移到$x=vt$，
回路面积$S(t)=\ell vt$，穿过的磁通量
\[
\Phi(t)=B_\perp(t)S(t)
=\frac12k\ell v t^2.
\]
故
\[
\boxed{|\mathcal E(t)|
=\left|\frac{d\Phi}{dt}\right|
=k\ell vt}.
\]
这一总感应电动势可分成磁场自身变化贡献
$k\ell vt/2$与杆运动扩大面积贡献$k\ell vt/2$。
\crosscheck 两项大小\emph{相等}，
漏掉动生或漏掉感生电动势都将得到一半结果。
若从磁场法线正侧观察，正法向通量增长，
由楞次定律感应电流沿\emph{顺时针}方向。

\sourceproblem{9.3.6}{239--240}{239--240}
\problemtext 一边长$a$的方形导线框以恒速$v$向右通过宽度$3a$的均匀磁场区，
磁场方向指向纸内。
$t=0$时方框的右侧前缘刚进入磁场。
以逆时针电流为正，写出整个进出磁场过程中感应电动势的分段表达式。
\workings
在进入阶段$0<t<a/v$，框内与磁场重叠的面积
$S(t)=avt$，指向纸内的磁通量增长，
感应电流沿逆时针，故
$\mathcal E=+Bav$。
当$a/v<t<3a/v$，整个线框位于均匀磁场内，
$S=a^2$恒定，电动势为零。
离开阶段$3a/v<t<4a/v$，重叠面积
$S(t)=a(4a-vt)$，指向纸内的通量减小，
感应电流沿顺时针，故电动势为负。
因此
\[
\boxed{\mathcal E(t)=
\begin{cases}
0,&t<0,\\
+Bav,&0<t<a/v,\\
0,&a/v<t<3a/v,\\
-Bav,&3a/v<t<4a/v,\\
0,&t>4a/v.
\end{cases}}
\]
\crosscheck 进场和离场脉冲持续时间均为$a/v$，
中间零电动势平台宽$2a/v$，
两脉冲积分的有向总和为零，符合初末磁通量同为零。
